% =====================================================================
%  GUIDE D'UTILISATION — LEAKO (Leaks Survey)
%  Guide simple et pratique pour un nouvel utilisateur
%  Compiler avec : pdflatex (ou Overleaf)
% =====================================================================
\documentclass[11pt,a4paper]{article}

% ------------------------- PACKAGES -------------------------
\usepackage[utf8]{inputenc}
\usepackage[T1]{fontenc}
\usepackage[french]{babel}
\usepackage{graphicx}
\usepackage{geometry}
\usepackage{xcolor}
\usepackage{titlesec}
\usepackage{fancyhdr}
\usepackage{hyperref}
\usepackage{enumitem}
\usepackage{booktabs}
\usepackage{float}
\usepackage{helvet}
\usepackage{setspace}

% ------------------------- MISE EN PAGE -------------------------
\geometry{top=2cm, bottom=2cm, left=2.5cm, right=2.5cm}
\renewcommand{\familydefault}{\sfdefault}
\setstretch{1.15}

% ------------------------- COULEURS OCP -------------------------
\definecolor{ocpvert}{RGB}{0, 122, 92}
\definecolor{ocpvertfonce}{RGB}{0, 82, 62}
\definecolor{gris}{RGB}{90, 90, 90}

% ------------------------- EN-TÊTE / PIED DE PAGE -------------------------
\pagestyle{fancy}
\fancyhf{}
\fancyhead[L]{\small\textcolor{ocpvert}{\textbf{LEAKO — Guide d'utilisation}}}
\fancyhead[R]{\small\textcolor{gris}{OCP}}
\fancyfoot[C]{\thepage}
\renewcommand{\headrulewidth}{0.4pt}
\renewcommand{\headrule}{\hbox to\headwidth{\color{ocpvert}\leaders\hrule height \headrulewidth\hfill}}

% ------------------------- TITRES -------------------------
\titleformat{\section}
  {\normalfont\Large\bfseries\color{ocpvertfonce}}
  {\thesection}{1em}{}
\titleformat{\subsection}
  {\normalfont\large\bfseries\color{ocpvert}}
  {\thesubsection}{1em}{}

% ------------------------- LIENS -------------------------
\hypersetup{
  colorlinks=true,
  linkcolor=ocpvertfonce,
  urlcolor=ocpvert,
}

% ------------------------- DOCUMENT -------------------------
\begin{document}

% =====================================================================
%  TITRE
% =====================================================================
\begin{center}
{\Huge\bfseries\textcolor{ocpvertfonce}{LEAKO}}\\[0.3cm]
{\Large\bfseries Guide d'utilisation rapide}\\[0.2cm]
{\large Gestion et inspection des fuites de vapeur}\\[0.5cm]
\rule{0.5\textwidth}{1pt}
\end{center}

\vspace{0.5cm}

% =====================================================================
%  1. INSTALLATION DE L'APPLICATION
% =====================================================================
\section{Installer l'application mobile}

\begin{enumerate}[leftmargin=2em]
    \item Ouvrez le \textbf{dashboard web} dans votre navigateur : \url{https://leako.quaisse.me}
    \item Allez sur la page \textbf{Application mobile}.
    \item \textbf{Scannez le QR code} avec votre téléphone pour télécharger l'application (fichier APK).
    \item Installez l'application sur votre téléphone Android.
\end{enumerate}

\subsection{Si vous n'avez pas encore l'application}
\begin{itemize}
    \item Le QR code se trouve sur la page \textbf{Application mobile} du dashboard.
    \item Scannez-le avec l'appareil photo de votre téléphone, puis téléchargez et installez l'APK.
\end{itemize}

% =====================================================================
%  2. CRÉATION DE COMPTE
% =====================================================================
\section{Créer un compte}

\begin{enumerate}[leftmargin=2em]
    \item Ouvrez l'application \textbf{LEAKO} sur votre téléphone.
    \item Sur l'écran de connexion, appuyez sur \textbf{Créer un compte}.
    \item Renseignez vos informations : \textbf{nom, prénom, email, mot de passe}.
    \item Validez. Vous êtes maintenant connecté.
\end{enumerate}

% =====================================================================
%  3. CRÉATION D'UN PROJET
% =====================================================================
\section{Créer un projet}

\begin{enumerate}[leftmargin=2em]
    \item Connectez-vous à l'application.
    \item Ouvrez le menu \textbf{Projets}.
    \item Appuyez sur \textbf{Nouveau projet}.
    \item Donnez un \textbf{nom} et une \textbf{description} au projet.
    \item Validez. Le projet est créé.
\end{enumerate}

% =====================================================================
%  4. INVITATION DE MEMBRES
% =====================================================================
\section{Inviter des membres}

\begin{enumerate}[leftmargin=2em]
    \item Ouvrez votre \textbf{projet}.
    \item Appuyez sur \textbf{Inviter un membre}.
    \item Saisissez l'\textbf{email} de la personne à inviter.
    \item La personne reçoit une invitation et doit l'\textbf{accepter}.
\end{enumerate}

\subsection{Important}
\begin{itemize}
    \item \textbf{Seuls les membres d'un projet} peuvent voir les données (fuites, campagnes) de ce projet.
    \item Une personne \textbf{hors projet} ne peut \textbf{pas} voir les données.
    \item Pour accéder aux données, il faut d'abord être \textbf{invité} et \textbf{accepter} l'invitation.
\end{itemize}

% =====================================================================
%  5. CRÉATION D'UNE CAMPAGNE
% =====================================================================
\section{Créer une campagne d'inspection}

\begin{enumerate}[leftmargin=2em]
    \item Ouvrez le menu \textbf{Campagnes}.
    \item Appuyez sur \textbf{Nouvelle campagne}.
    \item Renseignez le \textbf{nom} et la \textbf{zone} d'inspection.
    \item Validez. La campagne est créée.
\end{enumerate}

% =====================================================================
%  6. SIGNALEMENT D'UNE FUITE
% =====================================================================
\section{Signaler une fuite}

\begin{enumerate}[leftmargin=2em]
    \item Ouvrez le menu \textbf{Fuites}.
    \item Appuyez sur \textbf{Nouvelle fuite}.
    \item Renseignez les informations : \textbf{tag, type de vapeur, pression, diamètre}.
    \item \textbf{Capturez la position GPS} (bouton GPS).
    \item \textbf{Ajoutez des photos} et/ou une vidéo de la fuite.
    \item Validez. La fuite est enregistrée et le \textbf{coût des pertes} est calculé automatiquement.
\end{enumerate}

% =====================================================================
%  7. SUIVI ET ANALYSE
% =====================================================================
\section{Suivre et analyser}

\begin{itemize}
    \item \textbf{Suivi} : changez le statut d'une fuite (À réparer, En cours, Réparée, Annulée).
    \item \textbf{Analyse IA} : lancez l'analyse d'une photo/vidéo pour détecter automatiquement la fuite.
    \item \textbf{Rapports} : consultez les métriques et exportez les rapports en PDF.
    \item \textbf{Dashboard web} : supervisez toutes les données depuis le navigateur.
\end{itemize}

% =====================================================================
%  RÉCAPITULATIF
% =====================================================================
\section{Récapitulatif des étapes}

\begin{table}[H]
\centering
\begin{tabular}{@{}ll@{}}
\toprule
\textbf{Étape} & \textbf{Action} \\
\midrule
1 & Scanner le QR code et installer l'application \\
2 & Créer un compte \\
3 & Créer un projet \\
4 & Inviter des membres (ils doivent accepter) \\
5 & Créer une campagne d'inspection \\
6 & Signaler les fuites (GPS + photos) \\
7 & Suivre et analyser (statuts, IA, rapports) \\
\bottomrule
\end{tabular}
\caption{Les étapes clés pour démarrer}
\end{table}

\vspace{0.5cm}

\begin{center}
\rule{0.5\textwidth}{1pt}\\[0.3cm]
{\small\textcolor{gris}{Pour toute question, contactez votre administrateur.}}
\end{center}

\end{document}
